\documentclass[10pt,aspectratio=169]{beamer}

%-------------------------------------------------
% Theme and packages
%-------------------------------------------------
\usetheme{Madrid}
\usecolortheme{seahorse}
\useoutertheme{infolines}
\useinnertheme{rounded}

\usepackage[utf8]{inputenc}
\usepackage[T1]{fontenc}
\usepackage{amsmath,amssymb}
\usepackage{graphicx}
\usepackage{xcolor}
\usepackage{booktabs}
\usepackage{makecell}
\usepackage{hyperref}

% Slightly tighter list spacing
\setbeamertemplate{itemize items}[circle]
\setbeamertemplate{itemize subitem}[triangle]
\setbeamertemplate{navigation symbols}{}

% Highlight box for reviewer quotes
\setbeamercolor{quotebox}{bg=blue!8,fg=black}
\newcommand{\reviewerquote}[1]{%
  \begin{beamercolorbox}[sep=6pt,rounded=true,shadow=true]{quotebox}%
  \itshape #1%
  \end{beamercolorbox}%
}

%-------------------------------------------------
% Title information
%-------------------------------------------------
\title[SFF Evaluation -- Interpretation]{SFF discussions}
\author{Agenda June 12, SFF phase 2}
\institute[UiO]{University of Oslo}
\date{\today}

\begin{document}

%-------------------------------------------------
\begin{frame}
  \titlepage
\end{frame}

%-------------------------------------------------
\begin{frame}{Outline}
  \tableofcontents
\end{frame}


%=================================================
\section{Agenda}
%=================================================
\begin{frame}{June 12 agenda}
  \begin{itemize}
    \item Short summary from previous meeting (June 5)
    \item Localization
    \item Budget (Lasse)
    \item Planning of next meetings and writing
    \item Discussion of evaluation (if time)
    \item AOB
  \end{itemize}
\end{frame}


%=================================================
\section{Localization}
%=================================================
\begin{frame}{Starting discussion about localization}
  \begin{itemize}
    \item Lead suggestion: Justin (agreed) and Thomas
    \item Rough estimate on how many we could be
    \item Need to get started as soon as possible, see also forms from the MN-college (mail from Ingse N June 8)
  \end{itemize}
\end{frame}


\begin{frame}{Potential students, CS Physics and Chemistry, approx 14-16}
\begin{table}[htbp]
\centering
\begin{tabular}{lrrrrrr}
\hline
\textbf{Retning} & \textbf{EU-søkere} & \textbf{EU-1.pri} & \textbf{EU-tilbud} & \textbf{EU-ja} & \textbf{Norden-søkere} & \textbf{Norden-1.pri} \\
\hline
AMRA     & 16 & 3  & 0 & 0 & 62 & 16 \\
Astro    & 10 & 3  & 0 & 0 & 17 & 7  \\
Bio      & 21 & 10 & 2 & 1 & 38 & 14 \\
Kjemi    & 2  & 1  & 1 & 1 & 11 & 6  \\
Geo      & 8  & 3  & 1 & 1 & 25 & 10 \\
IFI      & 12 & 1  & 0 & 0 & 50 & 14 \\
Mekanikk & 5  & 0  & 0 & 0 & 13 & 4  \\
Fysikk   & 9  & 2  & 3 & 2 & 63 & 19 \\
\hline
\textbf{Total} & 83 & 23 & 7 & 5 & 279 & 90 \\
\hline
\end{tabular}
\end{table}
\end{frame}

\begin{frame}{Quantum, Mena and other, possibly 15-20}
\begin{table}[htbp]
\centering
\small
\begin{tabular}{lrrrr}
\hline
\textbf{Studieprogram} &
\textbf{Opptaksramme} &
\textbf{Søknader} &
\textbf{Søkere} &
\textbf{1.prioritet} \\
\hline
Data Science                                     & 20  & 312  & 312 & 148 \\
Elektronikk, informatikk og teknologi            & 30  & 114  & 88  & 48  \\
Entreprenørskap og innovasjonsledelse            & 25  & 337  & 337 & 140 \\
Fornybare energisystemer                         & 30  & 119  & 119 & 54  \\
Fysikk                                           & 30  & 180  & 120 & 54  \\
Kjemi                                            & 25  & 204  & 105 & 74  \\
Kvantevitenskap og kvanteteknologi               & 15  & 47   & 47  & 23  \\
Matematikk                                       & 25  & 98   & 79  & 39  \\
Materialvitenskap for energi- og nanoteknologi   & 15  & 39   & 39  & 19  \\
Nukleærteknologi                                 & 15  & 47   & 47  & 20  \\
Stokastisk modellering, statistikk og risikoanalyse & 15 & 219 & 188 & 74 \\
\hline
\end{tabular}
\end{table}
\end{frame}

\begin{frame}{PIs and associated researchers (1/2)}
\begin{table}[h]
\centering
\scriptsize
\setlength{\tabcolsep}{3.5pt}
\renewcommand{\arraystretch}{1.3}
\begin{tabular}{@{}>{\raggedright\arraybackslash}p{2.7cm}>{\raggedright\arraybackslash}p{1.9cm}llcccc>{\raggedright\arraybackslash}p{2.3cm}c@{}}
\toprule
\textbf{Name} & \textbf{Title} & \textbf{Aff.} & \textbf{Role} & \textbf{RT1} & \textbf{RT2} & \textbf{RT3} & \textbf{RT4} & \textbf{Ext.\ funding} & \textbf{m/f}\\
\midrule
\multicolumn{10}{@{}l}{\textit{Centre directors}}\\
Morten Hjorth-Jensen      & Professor      & Phys., UiO & Dir.$^{a}$/PI & X & X & X  & X & NSF, DOE, RCN, QuantERA          & m\\
Marianne E.\ Bathen       & Assoc.\ Prof.\ & Phys., UiO & Dir.$^{b}$/PI & X &   & X & X & YRT$^{c}$, Equinor     & f\\
\midrule
\multicolumn{10}{@{}l}{\textit{Principal investigators}}\\
Thomas Bondo Pedersen     & Professor      & Chem., UiO & PI            & X & X &   & X & RP$^{d}$               & m\\
Simen Kvaal               & Researcher     & Chem., UiO & PI            & X & X &   & X & ERC-StG, CAS$^{e}$     & m\\
Michal Repisk\'y          & Sr.\ Researcher& Chem., UiO & PI            & X & X &   & X & MSCA-IF$^{f}$          & m\\
Lasse Vines               & Professor      & Phys., UiO & PI            & X &   & X &   & Toppforsk (RCN)        & m\\
Justin W.\ Wells          & Professor      & Phys., UiO & PI            & X &   & X &   & RCN                    & m\\
\bottomrule
\end{tabular}
\end{table}
\end{frame}

\begin{frame}{PIs and associated researchers (2/2)}
\begin{table}[h]
\centering
\scriptsize
\setlength{\tabcolsep}{3.5pt}
\renewcommand{\arraystretch}{1.3}
\begin{tabular}{@{}>{\raggedright\arraybackslash}p{2.7cm}>{\raggedright\arraybackslash}p{1.9cm}llcccc>{\raggedright\arraybackslash}p{2.3cm}c@{}}
\toprule
\textbf{Name} & \textbf{Title} & \textbf{Aff.} & \textbf{Role} & \textbf{RT1} & \textbf{RT2} & \textbf{RT3} & \textbf{RT4} & \textbf{Ext.\ funding} & \textbf{m/f}\\
\midrule
\multicolumn{10}{@{}l}{\textit{Associated investigators (AI)}}\\
Susanne Viefers       & Professor            & Phys., UiO     & AI & X & X &   &   &                & f\\
Francesco P.\ Massel  & Professor            & Phys., UiO/USN & AI & X & X & X & X & RCN, QuantERA  & m\\
Ainara Nova           & Assoc.\ Prof.\       & Chem., UiO     & AI &   &   &   &   &                & f\\
Michele Cascella      & Professor            & Chem., UiO     & AI &   &   &   & X &                & m\\
David Balcells        & Sr.\ Researcher      & Chem., UiO     & AI &   &   &   & X &                & m\\
Erik Tellgren         & Sr.\ Researcher      & Chem., UiO     & AI & X & X &   & X &                & m\\
Antoine Camper        & Sr.\ Researcher      & Phys., UiO     & AI & X & X &   &   &                & m\\
Simon Cooil           & Sr.\ Researcher      & Phys., UiO     & AI &   &   & X &   &                & m\\
David Gongorra        & Sr.\ Researcher      & Phys., UiO     & AI &   &   & X &   &                & m\\
Cecilie Glittum       & Post-doctoral fellow & Phys., UiO     & AI & X & X &   & X &                & f\\
Gunnar F.\ Lange      & Post-doctoral fellow & Phys., UiO     & AI & X & X &   & X &                & m\\
\bottomrule
\end{tabular}
\end{table}
\end{frame}


\begin{frame}{Total estimate}
  \begin{itemize}
  \item PIs, associated researchers and guests: 20
  \item Master of Science students: approx 30
  \item PhDs and PDs: estimate 15-20 at any time
  \item Total people: 65-70
  \end{itemize}
\end{frame}

%-------------------------------------------------
% Space-requirement form (translated from the UiO 2026
% "Skjema for Arealbehov", page 1, second table)
%-------------------------------------------------
\begin{frame}{Space requirements for a co-located centre}
  \small
  Form to be completed for a co-located SFF VI centre (2028--2038).\\[-0.3em]
  {\footnotesize\itshape Translated from the UiO 2026 ``Skjema for Arealbehov''.}
  \vfill
  \begin{table}[h]
  \centering
  \scriptsize
  \setlength{\tabcolsep}{4pt}
  \renewcommand{\arraystretch}{1.9}
  \begin{tabular}{@{}>{\raggedright\arraybackslash}p{3.0cm}cc>{\raggedright\arraybackslash}p{1.9cm}cc@{}}
  \toprule
  \textbf{Space category} &
  \makecell{\textbf{Area (m\textsuperscript{2})}\\\textbf{dry rooms}} &
  \makecell{\textbf{Area (m\textsuperscript{2})}\\\textbf{wet rooms}} &
  \textbf{Possible physical location} &
  \makecell{\textbf{Assumed}\\\textbf{conversion}\\\textbf{need (m\textsuperscript{2})}} &
  \makecell{\textbf{Estimated}\\\textbf{cost (NOK)}\textsuperscript{*}} \\
  \midrule
  Space in the institute's premises          & & & & & \\
  Space in the faculty's premises            & & & & & \\
  Space in UiO premises outside inst./fac.   & & & & & \\
  Space in external premises                 & & & & & \\
  Need for new space                         & & & & & \\
  \textbf{Total}                             & & & & & \\
  \bottomrule
  \end{tabular}
  \end{table}
  \vfill
  {\scriptsize
  \textsuperscript{*}\,Conversion costs are estimated at NOK 15\,000/m\textsuperscript{2} for dry area and
  NOK 25\,000--30\,000/m\textsuperscript{2} for wet area; the estimates already include relatively large rebuilds.}
\end{frame}

%-------------------------------------------------
\begin{frame}{Space requirements --- notes on rent}
  \begin{block}{Internal rent model}
    The internal-rent model applies: the MN-college is charged for the space it occupies
    according to the prevailing rates.
  \end{block}
  \begin{block}{External rent}
    External rent varies and must be requested from the relevant external owner. Price variations
    here are large and are clarified/approved by each individual faculty before any lease
    agreement is signed.
  \end{block}
\end{frame}

\section{Budget}

\begin{frame}{Budget --- preliminary cost breakdown}
  {\footnotesize\itshape Preliminary figures (Lasse); amounts in NOK over the 10-year centre period.}
  \vfill
  \footnotesize
  \setlength{\tabcolsep}{6pt}
  \renewcommand{\arraystretch}{1.15}
  \begin{columns}[T]
    \column{0.5\textwidth}
    \begin{tabular}{@{}lr@{}}
    \toprule
    \textbf{Item} & \textbf{Cost} \\
    \midrule
    9 PhDs (NFR funded)          & 48\,562\,500 \\
    7 PhDs (KD)                  & 0 \\
    5 postdocs (NFR funded)      & 26\,970\,500 \\
    Michal (10 yr)               & 18\,003\,000 \\
    Researcher (7 yr)            & 12\,565\,000 \\
    Centre Director (NFR funded) & 8\,286\,000 \\
    Administration (10 yr)       & 14\,807\,688 \\
    Sigma2 CPU (10M/yr)          & 4\,000\,000 \\
    \bottomrule
    \end{tabular}

    \column{0.5\textwidth}
    \begin{tabular}{@{}lr@{}}
    \toprule
    \textbf{Item} & \textbf{Cost} \\
    \midrule
    Sigma2 GPU (100k/yr)             & 5\,600\,000 \\
    Leiested (lab access)            & 3\,675\,000 \\
    Annual schools and workshops     & 7\,500\,000 \\
    Travel                           & 17\,500\,000 \\
    Centre running costs             & 2\,770\,000 \\
    Strategic funds, Faculty (drift) & 8\,250\,000 \\
    Storage (10 TB, 10 yr)           & 50\,000 \\
    \bottomrule
    \end{tabular}
  \end{columns}
  \vfill
  \begin{center}
    \normalsize\textbf{Total: 178\,539\,688 NOK}
  \end{center}
\end{frame}

%-------------------------------------------------
\begin{frame}{Budget --- income and key figures}
  \begin{columns}[T]
    \column{0.52\textwidth}
    \begin{block}{Income}
      \begin{itemize}
        \item 160 MNOK from NFR
        \item 20 MNOK from MN
      \end{itemize}
      \vspace{0.2em}
      \centering\textbf{Total income: 180 MNOK}
    \end{block}
    \begin{block}{Cost vs.\ income}
      Total cost $\approx$ 178.5 MNOK against 180 MNOK income.
    \end{block}

    \column{0.48\textwidth}
    \begin{block}{Key figures}
      \begin{itemize}
        \item Nettobidrag (net contribution): 61\%
        \item Overhead / indirekte kostnader (indirect costs): 48\%
        \item Eksternt finansiert overhead (ext.-financed overhead): 36\%
        \item Egenandel (own contribution): 39\%
      \end{itemize}
    \end{block}
  \end{columns}
  \vfill
  \begin{block}{Note on external funding}
    All external funding for the PIs has been removed, except the 50\,\% (compulsory)
    external funding for the Centre Director.
  \end{block}
\end{frame}


\section{Planning of meetings}

\begin{frame}{When and plans for writing}
  \begin{itemize}
  \item When and how often (zoom only most likely)
  \item First tentative draft to Innovayt week of July 8-13
  \end{itemize}
\end{frame}



%=================================================
\section{RCN feedback}
%=================================================
\begin{frame}{Overall impression of the Phase 1 evaluation}
  \begin{itemize}
    \item The review says this  was a \textbf{very strong} Phase 1 proposal.
    \item All four criteria received the mark \textbf{6 (Excellent)} on the 1--7 scale. Is this scale correct?
    \item On the Research Council scale, mark 6 means the proposal
          \emph{``successfully addresses essentially all aspects of the criterion and has only minor shortcomings.''}
    \item The path to a competitive Phase 2 application is therefore essentially about
          \textbf{closing the gap from 6 to 7 (Outstanding)} by addressing the specific
          \emph{minor shortcomings} hinted at in each section.
  \end{itemize}
  \vfill
  \begin{block}{Structure of these notes}
    \begin{enumerate}
      \item A less in-depth interpretation of the feedback.
      \item A discussion of explicit things we can address in the full proposal.
    \end{enumerate}
  \end{block}
\end{frame}

%=================================================
\section{Part 1: A less in-depth interpretation}
%=================================================

\begin{frame}{Excellence -- Quality of R\&D activities}
  This section contains the most concrete criticism. Two distinct points are flagged.
  \vfill
  \begin{block}{1. Research questions need sharper articulation}
    Each of the four tracks should open with a \emph{small set of falsifiable scientific questions},
    rather than thematic descriptions.
  \end{block}
  \begin{block}{2. Machine-learning methods need more detail}
    The ML-guided control component is the methodological area the panel found least convincing.
    The full application should include:
    \begin{itemize}
      \item concrete algorithms and model architectures,
      \item training and data strategies,
      \item validation protocols,
      \item the interface between ML and the physics/chemistry tracks.
    \end{itemize}
    A dedicated methods section for ML-guided quantum control is likely warranted.
  \end{block}
\end{frame}

%-------------------------------------------------
\begin{frame}{Impact}
  Two implicit concerns stand out.
  \vfill
  \begin{block}{Scale}
    The panel notes that the center is
    \reviewerquote{relatively small compared with major international centers}
    \vspace{0.3em}
    This is the only comparative remark in the whole report -- and it is not a compliment.
    Phase 2 should either
    \begin{itemize}
      \item grow the scope and team to match international peers, or
      \item explicitly argue why the chosen size delivers \emph{more impact per krone} than larger competitors.
    \end{itemize}
  \end{block}
  \begin{block}{Generic impact narrative}
    The current text stays at the level of \emph{application areas} without committing to specific deliverables,
    pathways, or partners. Phase 2 requires explicit dissemination, exploitation, and open-science
    plans, plus communication and engagement with different target audiences.
  \end{block}
\end{frame}

%-------------------------------------------------
\begin{frame}{Implementation}
  \reviewerquote{some uncertainty exists regarding continuity and transition of leadership}
  \vfill
  Succession uncertainty is one of the recurring reasons SFF panels downgrade otherwise excellent centers.
  \vfill
  \begin{block}{Phase 2 should include a written succession plan}
    \begin{itemize}
      \item Who can step in as director if needed?
      \item How are deputy and co-director roles structured?
      \item How is leadership of each of the four tracks backed up?
      \item How does the center survive the departure of a key PI?
    \end{itemize}
  \end{block}
  \begin{block}{Additional Phase 2 requirements (not yet present)}
    \begin{itemize}
      \item An explicit training environment beyond what individual groups offer.
      \item Appropriate management and governance structures.
      \item A description of the physical organisation supporting cross-track collaboration.
      \item If gender-imbalanced: concrete plans to qualify under-represented-gender researchers
            for senior positions.
    \end{itemize}
  \end{block}
\end{frame}

%=================================================
\section{Part 2: Explicit things we can address}
%=================================================

\begin{frame}{How to read the comments}
  Our goal is to prepare a much larger and more competitive Phase 2 SFF proposal.
  \vfill
  \begin{itemize}
    \item We should read the comments not as criticism, but as indications of where reviewers want
          \textbf{additional depth, precision, and evidence}.
    \item I think that for  highly competitive SFF evaluations, proposals with uniformly excellent scores are
          distinguished by how convincingly they address \emph{these} details.
  \end{itemize}
\end{frame}

%-------------------------------------------------
\begin{frame}{1. The clearest weakness: explicit research questions}
  \reviewerquote{The proposal would benefit from clearer articulation of the specific research questions within each track.}
  \vfill
  The reviewers liked the four-track structure but wanted more than broad themes.
  \vfill
  \begin{columns}[T]
    \column{0.45\textwidth}
    \textbf{Instead of:}
    \begin{itemize}
      \item Quantum control of many-body systems
      \item Machine learning for quantum dynamics
    \end{itemize}

    \column{0.55\textwidth}
    \textbf{Formulate questions such as:}
    \begin{itemize}\small
      \item What are the fundamental limits on controllability in interacting many-body systems?
      \item Can agent-based AI (with reinforcement learning) discover control protocols beyond Pontryagin-optimal solutions?
      \item Can quantum Fisher information provide universal bounds on control fidelity?
      \item Can generative models compress experimentally relevant quantum states?
    \end{itemize}
  \end{columns}
  \vfill
  \begin{center}
    \emph{The best SFF proposals resemble a collection of major open problems, not activities.}
  \end{center}
\end{frame}

%-------------------------------------------------
\begin{frame}{2. Machine learning needs greater methodological precision}
  \reviewerquote{more detailed descriptions of certain methods -- particularly those related to machine-learning-guided control.}
  \vfill
  The panel accepted that ML is important, but were not yet convinced it is sufficiently specified.
  \vfill
  \begin{columns}[T]
    \column{0.45\textwidth}
    \textbf{Name explicit ML methodologies:}
    \begin{itemize}\small
      \item Agent-based AI and Reinforcement learning
      \item Physics-informed neural networks
      \item Bayesian optimization
      \item Diffusion models
      \item Transformer architectures
      \item Generative quantum state models
      \item Graph neural networks
    \end{itemize}

    \column{0.5\textwidth}
    \textbf{For each method, explain:}
    \begin{itemize}\small
      \item why it is needed,
      \item what scientific question it addresses,
      \item why existing methods are insufficient.
    \end{itemize}
  \end{columns}
\end{frame}

%-------------------------------------------------
\begin{frame}{2. ML as a research theme, not just a tool}
  A common weakness in AI-related proposals is that ML appears as a \emph{tool} rather than a
  \emph{research theme}.
  \vfill
  \begin{block}{The proposal should argue}
    New machine-learning methodologies will emerge \textbf{from quantum control challenges themselves}.
  \end{block}
  \vfill
  \begin{block}{Examples of methodological novelty}
    \begin{itemize}
      \item Reinforcement learning for quantum control under uncertainty.
      \item Scientific foundation models for quantum experiments.
      \item Physics-informed generative models.
      \item Hybrid quantum-classical learning algorithms.
    \end{itemize}
  \end{block}
\end{frame}

%-------------------------------------------------
\begin{frame}{3. Clarify scientific ambition versus feasibility}
  The reviewers repeatedly praise the ambition, but Phase 2 should make the relationship
  between risk levels much more explicit.
  \vfill
  \begin{columns}[T]
    \column{0.5\textwidth}
    \textbf{Flagship challenges}
    \begin{itemize}
      \item Fault-tolerant quantum control.
      \item AI-designed quantum sensors.
      \item Autonomous quantum experiments.
    \end{itemize}

    \column{0.5\textwidth}
    \textbf{Fallback pathways}\\
    If a flagship goal fails:
    \begin{itemize}
      \item What is still learned?
      \item What publishable results remain?
    \end{itemize}
  \end{columns}
  \vfill
  \begin{center}
    SFF panels routinely ask: \emph{What happens if the most ambitious ideas fail?}
  \end{center}
\end{frame}

%-------------------------------------------------
\begin{frame}{4. Leadership succession and long-term sustainability}
  \reviewerquote{some uncertainty exists regarding continuity and transition of leadership}
  \vfill
  The panel implicitly asked:
  \begin{itemize}
    \item How dependent is the center on a single person?
    \item What happens after 5--10 years?
    \item How are younger leaders being developed?
  \end{itemize}
  \vfill
  \begin{block}{Explicit leadership pipeline}
    \begin{itemize}
      \item Deputy director(s).
      \item Future work-package leaders.
      \item Early-career leadership positions.
      \item Documented succession plans.
    \end{itemize}
  \end{block}
\end{frame}

%-------------------------------------------------
\begin{frame}{5. Demonstrate stronger center-level synergy}
  The next proposal must repeatedly answer:
  \begin{center}
    \large \emph{Why does this need to be a Center of Excellence, rather than a collection of excellent groups?}
  \end{center}
  \vfill
  \begin{block}{Cross-track flagship projects}
    \begin{itemize}
      \item Quantum sensing $+$ AI.
      \item Quantum materials $+$ machine learning.
      \item Quantum chemistry $+$ control theory.
      \item Experimental platforms $+$ digital twins.
    \end{itemize}
  \end{block}
  \vfill
  \begin{center}
    Each major breakthrough should \textbf{require} collaboration between several tracks.
  \end{center}
\end{frame}

%-------------------------------------------------
\begin{frame}{6. Strengthen international positioning}
  \reviewerquote{relatively small compared with major international centers}
  \vfill
  This reveals a comparison class. For Phase 2, compare explicitly to:
  \vfill
  \begin{columns}[T]
    \column{0.5\textwidth}
    \begin{itemize}
      \item QuICS (Maryland)
      \item IQIM (Caltech)
      \item IQC Waterloo
    \end{itemize}

    \column{0.5\textwidth}
    \begin{itemize}
      \item CQT Singapore
      \item Munich Quantum Valley
      \item Quantum Delta NL
    \end{itemize}
  \end{columns}
  \vfill
  \begin{block}{Then explain}
    \begin{itemize}
      \item Where Norway can \emph{lead}.
      \item What niche is uniquely Norwegian.
      \item Why the center will become internationally indispensable.
    \end{itemize}
  \end{block}
\end{frame}

%-------------------------------------------------
\begin{frame}{7. Expand the impact narrative}
  The current impact section is strong but generic (secure communication, advanced materials,
  healthcare, clean energy). Phase 2 should make this much more concrete.
  \vfill
  \begin{columns}[T]
    \column{0.33\textwidth}
    \textbf{Scientific impact}
    \begin{itemize}\small
      \item New theory of quantum control.
      \item New algorithms.
      \item Open-source software.
    \end{itemize}

    \column{0.33\textwidth}
    \textbf{Technology impact}
    \begin{itemize}\small
      \item Quantum sensors.
      \item Quantum materials discovery.
      \item Quantum computing architectures.
    \end{itemize}

    \column{0.33\textwidth}
    \textbf{Human capital impact}
    \begin{itemize}\small
      \item Number of PhDs.
      \item Industrial internships.
      \item International mobility.
    \end{itemize}
  \end{columns}
\end{frame}

%-------------------------------------------------
\begin{frame}{8. What is missing entirely from the criticism?}
  Interestingly, there are \textbf{no} criticisms regarding:
  \begin{itemize}
    \item scientific quality,
    \item leadership competence,
    \item publication record,
    \item feasibility,
    \item interdisciplinarity.
  \end{itemize}
  \vfill
  \begin{center}
    \emph{This is very encouraging.} The panel is convinced the team can do the work.
  \end{center}
  \vfill
  \begin{block}{The next proposal should focus less on proving excellence and more on:}
    \begin{enumerate}
      \item Sharper scientific questions.
      \item More detailed methodology.
      \item Clearer ML strategy.
      \item Stronger center-level integration.
      \item Leadership succession.
      \item A stronger international vision.
    \end{enumerate}
  \end{block}
\end{frame}

%=================================================
\section{Overall interpretation}
%=================================================

\begin{frame}{Reading between the lines}
  \reviewerquote{We are convinced that this is an excellent team pursuing an important and
  potentially transformative topic. For the full SFF proposal we would like to see a more
  detailed scientific roadmap, especially regarding machine-learning-guided quantum control,
  clearer articulation of the key scientific questions, and stronger evidence that the center
  will function as more than the sum of its individual research groups.}
  \vfill
  \begin{itemize}
    \item A very favourable position going into a second-round SFF application.
    \item The reviewers did not identify fundamental weaknesses.
    \item They asked for \textbf{greater precision}, \textbf{sharper focus}, and a
          \textbf{more developed center strategy}.
  \end{itemize}
\end{frame}

%-------------------------------------------------
\begin{frame}{Summary: priorities for the Phase 2 proposal}
  \begin{enumerate}
    \item \textbf{Leadership continuity and succession plan} -- the only explicit Implementation concern.
    \item \textbf{Detailed methodology for ML-guided quantum control} -- dedicated methods section.
    \item \textbf{Sharpened research questions} per track -- falsifiable, ambitious, specific.
    \item \textbf{Scale and international positioning} -- either grow, or justify the chosen scale against named peer centers.
    \item \textbf{Center-level synergy} -- cross-track flagship projects that no single group could deliver.
    \item \textbf{Risk management} -- flagship goals with explicit fallback pathways.
    \item \textbf{Phase 2 requirements} -- dissemination, exploitation, open science, training environment, gender plans, governance.
  \end{enumerate}
\end{frame}


\end{document}


